\section{Embedding}\label{sec:embedding}

Tra l'indicizzazione e il recupero si colloca un componente condiviso dai due processi (Figura~\ref{fig:componenti}):
l'Embedding trasforma ogni frammento, al momento dell'indicizzazione, nella sua rappresentazione vettoriale, e la
stessa trasformazione viene applicata alla domanda di uno studente al momento del recupero, così che domanda e
frammenti siano confrontabili nello stesso spazio (Sezione~\ref{sec:rag}). Questa sezione descrive il modello usato,
cosa viene effettivamente trasformato in vettore e come viene conservato nel database vettoriale.

\subsection{Il modello: BGE-M3}

Il modello di embedding adottato è BGE-M3~\cite{chen2024bgem3}, già discusso nella Sezione~\ref{subsec:embedding}. La
scelta è motivata da una caratteristica specifica di questo modello più che dalla sola qualità delle sue
rappresentazioni: a differenza di un modello di embedding classico, che produce solo una rappresentazione densa, BGE-M3
ne produce tre nello stesso passaggio di calcolo --- una densa, una sparsa e una per token --- evitando così di dover addestrare, eseguire e
mantenere allineati modelli distinti per ciascun tipo di rappresentazione. Un secondo motivo è la finestra di
contesto del modello, fino a 8.192 token: più ampia del limite di 512 token scelto per un frammento, lascia margine a non troncare mai un frammento durante l'embedding, anche per quelli
risultanti da una fusione gerarchica di più frammenti originari.

A titolo di esempio, per il frammento della Tabella~\ref{tab:metadati} --- il materiale didattico del corso di
Tecniche di Consulenza nella Scuola sul \emph{Change Laboratory} --- BGE-M3 calcola, nello stesso passaggio, le tre
rappresentazioni seguenti:
\begin{itemize}
  \item una rappresentazione \textbf{densa} di 1.024 numeri in virgola mobile, le cui prime componenti sono, in
    questo caso, $-0{,}009,\ -0{,}036,\ 0{,}007,\ 0{,}022,\ -0{,}007,\ -0{,}028,\ \dots$;
  \item una rappresentazione \textbf{sparsa}, che assegna un peso ai soli token lessicalmente rilevanti del testo;
    la Tabella~\ref{tab:sparse-esempio} mette a confronto, per questo frammento, il peso di alcune parole chiave
    con quello, nullo, di un articolo e di una preposizione --- la rappresentazione sparsa distingue così le parole
    che portano significato da quelle puramente grammaticali;
  \item una rappresentazione \textbf{per token}, una matrice di 87 vettori da 1.024 numeri ciascuno --- uno per
    ogni token in cui il testo del frammento viene suddiviso --- usata dalla sola modalità di recupero con late
    interaction (Sezione~\ref{sec:rag}).
\end{itemize}

\begin{table}[H]
  \centering
  \small
  \begin{tabular}{@{} l r @{}}
    \toprule
    \textbf{Parola} & \textbf{Peso} \\
    \midrule
    \texttt{intervento} & 0,236 \\
    \texttt{consulenza} & 0,234 \\
    \texttt{Ciclo} & 0,148 \\
    \texttt{sviluppato} & 0,125 \\
    \midrule
    \texttt{il} & 0,000 \\
    \texttt{a} & 0,000 \\
    \bottomrule
  \end{tabular}
  \caption{Peso, nella rappresentazione sparsa del frammento \emph{Change Laboratory} della Tabella~\ref{tab:metadati},
  di quattro parole che portano significato contro quello di un articolo e di una preposizione: la rappresentazione
  sparsa assegna peso nullo alle seconde, pur comparendo entrambe nel testo.}
  \label{tab:sparse-esempio}
\end{table}

\subsection{Cosa viene trasformato in vettore}

Non tutto il testo associato a un frammento entra nella sua rappresentazione vettoriale. Il percorso delle
intestazioni della sezione di provenienza (Sezione~\ref{sec:chunking}) viene incluso, perché è spesso necessario a
rendere comprensibile un frammento altrimenti ambiguo; gli altri metadati --- corso, ambito, sezione, stato,
l'indirizzo della pagina di provenienza e gli altri campi della Tabella~\ref{tab:metadati} --- ne restano invece
esclusi, e intervengono solo più avanti, come filtri sul database vettoriale (Sezione~\ref{sec:ricerca-ibrida}), non
come parte del significato che l'embedding rappresenta: il nome del corso di un frammento, ad esempio, non deve
influenzare quanto quel frammento è semanticamente vicino a una domanda, un compito che spetta al filtro, non alla
rappresentazione vettoriale.

\subsection{Il database vettoriale: Qdrant}

Le rappresentazioni calcolate vengono salvate in Qdrant, il database vettoriale scelto per due capacità native che
evitano di doverle replicare lato applicazione. La prima è la fusione tra ricerca densa e sparsa con Reciprocal Rank
Fusion (Sezione~\ref{sec:rag}), eseguita da Qdrant stesso a partire da due vettori nominati distinti, \emph{dense} e
\emph{sparse}, associati allo stesso punto; un terzo vettore nominato, \emph{colbert}, multivettore con
confronto per massima similarità, ospita la rappresentazione per token usata dalla sola modalità di recupero con
late interaction. La seconda è l'indicizzazione dei campi del
payload --- stato, corso, ambito, tipo di laurea, sezione, insegnamento e l'identificativo del documento --- come
indici di tipo \emph{keyword}: permettono di filtrare i candidati per questi campi nello stesso passaggio della
ricerca vettoriale, invece di recuperare un insieme più ampio di candidati e scartare quelli con i metadati sbagliati
in un passaggio successivo.

Ogni punto della collezione conserva, oltre ai tre vettori, il testo del frammento e l'insieme completo dei suoi
metadati (Sezione~\ref{sec:metadati}) come payload.

Il Listato~\ref{lst:qdrant-point} riporta, abbreviato, il punto realmente salvato in Qdrant per il frammento
\emph{Change Laboratory} della Tabella~\ref{tab:metadati} e della Tabella~\ref{tab:sparse-esempio}: ogni vettore
nominato (\texttt{dense}, \texttt{sparse}, \texttt{colbert}) accanto al payload che lo accompagna, con i campi
indicizzati --- \texttt{ambito}, \texttt{corso}, \texttt{stato} e gli altri citati sopra --- pronti per il filtro
applicato in fase di ricerca (Sezione~\ref{sec:ricerca-ibrida}).

\begin{lstlisting}[caption={Punto Qdrant del frammento \emph{Change Laboratory} (identificativo abbreviato nei
valori numerici, per brevità).}, label={lst:qdrant-point}]
{
  "id": "f23618c3-3550-5de1-8684-c1a92cc6a187",
  "vector": {
    "dense":   [-0.009, -0.036, 0.007, 0.022, -0.007, -0.028, ...],  // 1024 numeri
    "sparse":  {"indices": [21488, 30011, 111636, ...],
                "values":  [0.320, 0.274, 0.236, ...]},              // ~60 coppie non nulle
    "colbert": [[...], [...], ..., [...]]                           // 87 vettori da 1024 numeri
  },
  "payload": {
    "id": "3e53833e-ffc7-5e84-ab82-9a13a6776148",
    "ambito": "corso",
    "tipo_laurea": "magistrale",
    "corso": "PsicologiaScolasticaComunita",
    "sezione": "corsi",
    "materia": "TecnicheDiConsulenzaNellaScuola",
    "headings": ["Change Laboratory (Change Lab)"],
    "stato": "vigente",
    "checksum": "8e837f4622a5b4e2",
    "text": "Metodo di intervento sviluppato da Engestrom, basato sull'Activity Theory. ..."
  }
}
\end{lstlisting}

\subsection{L'indice a grafo: HNSW}

Oltre a salvare vettori e payload, Qdrant costruisce, per il vettore denso --- e separatamente per quello
\emph{colbert} --- la struttura dati usata a tempo di ricerca per trovare i vicini più prossimi alla domanda senza
confrontarla con ogni singolo punto della collezione: un grafo secondo lo schema Hierarchical Navigable Small World
(HNSW), introdotto in astratto nella Sezione~\ref{subsec:embedding}. Ogni punto è un nodo, collegato a un piccolo
insieme di altri punti vicini nello spazio degli embedding; i nodi sono distribuiti su più livelli, con parametri
che per questa collezione sono quelli di default di Qdrant, non modificati: ogni nodo ha al più $m=16$ collegamenti
per livello e, in costruzione, ne valuta $\mathit{ef\_construct}=100$ candidati. Il livello più alto contiene solo
una piccola frazione dei punti, collegati tra loro a lungo raggio; il livello più basso li contiene tutti, con
collegamenti brevi verso i vicini effettivi (Figura~\ref{fig:hnsw}). Una ricerca parte da un punto di ingresso nel
livello più alto e scende un livello alla volta, spostandosi a ogni passo verso il nodo già esplorato più vicino
alla domanda, fino a raggiungere il livello più basso, dove l'esplorazione dell'intorno del nodo raggiunto
restituisce i candidati da passare al riordinamento (Sezione~\ref{sec:corso-priorita}). È questo grafo a rendere praticabile la ricerca densa dello stadio di ricerca ibrida
(Sezione~\ref{sec:ricerca-ibrida}) con un tempo che cresce in modo logaritmico col numero di frammenti indicizzati
invece che linearmente.

\begin{figure}[H]
  \centering
  \begin{tikzpicture}[
      node/.style={draw, circle, minimum size=0.5cm, inner sep=0pt, black!60, fill=white},
      entry/.style={node, fill=blue!18, draw=blue!60!black, thick},
      hit/.style={node, fill=orange!35, draw=orange!70!black, thick},
      lyedge/.style={black!35, thin},
      samecol/.style={black!25, thin, dotted},
      srch/.style={-{Stealth[length=2.2mm]}, thick, red!65!black},
      lylabel/.style={font=\scriptsize, black!70, anchor=east},
    ]
    \foreach \x in {0,4,8} \draw[samecol] (\x,0) -- (\x,4.4);

    \foreach \x/\name in {0/p0,1/p1,2/p2,3/p3,4/p4,5/p5,6/p6,7/p7,8/p8}
      \node[node] (\name) at (\x,0) {};
    \draw[lyedge] (p0)--(p1)--(p2)--(p3)--(p4)--(p5)--(p6)--(p7)--(p8);
    \node[hit] at (p7) {};
    \node[lylabel] at (-1.1,0) {Livello 0};

    \foreach \x/\name in {0/m0,2/m2,4/m4,6/m6,8/m8}
      \node[node] (\name) at (\x,2.2) {};
    \draw[lyedge] (m0)--(m2)--(m4)--(m6)--(m8);
    \node[lylabel] at (-1.1,2.2) {Livello 1};

    \foreach \x/\name in {0/t0,4/t4,8/t8}
      \node[node] (\name) at (\x,4.4) {};
    \draw[lyedge] (t0)--(t4)--(t8);
    \node[entry] at (t0) {};
    \node[lylabel] at (-1.1,4.4) {Livello 2};

    \draw[srch] (t0) -- (t4);
    \draw[srch] (t4) -- (m4);
    \draw[srch] (m4) -- (m6);
    \draw[srch] (m6) -- (p6);
    \draw[srch] (p6) -- (p7);

    \node[font=\scriptsize, blue!50!black, above=2pt of t0] {ingresso};
    \node[font=\scriptsize, orange!70!black, below=2pt of p7] {vicino trovato};
  \end{tikzpicture}
  \caption{Ricerca del vicino più prossimo nel grafo HNSW (schematico): dal punto di ingresso nel livello più rado
  (Livello 2), la ricerca scende un livello alla volta muovendosi verso il nodo più vicino alla domanda a ogni
  passo, fino a raggiungere il livello più fitto (Livello 0, tutti i punti) ed esplorarne l'intorno. Le linee
  tratteggiate verticali collegano lo stesso punto presente su più livelli.}
  \label{fig:hnsw}
\end{figure}
